\section{Testing Strategy and Methodology}

Validation addresses algorithmic accuracy, numerical precision, system integration, and performance. The test suite follows a three-tier hierarchy: unit tests (component isolation), integration tests (component interactions), and validation tests (analytical comparison).

\subsection{Test Case Design}

\textbf{Boundary cases} probe edge conditions: minimal/maximal antenna counts, FFT sizes from 64 to 4096, and parameter extremes to expose off-by-one errors and buffer overflows.

\textbf{Known signal cases} use synthetic inputs with analytical solutions (pure tones, geometric delays) for precise verification.

\textbf{Realistic scenarios} simulate observing conditions: multiple sources, noise-dominated signals, overlapping integrations.

\textbf{Regression tests} prevent bug recurrence by capturing fixed failure conditions.

\subsection{Automated Testing Infrastructure}

Pytest discovers and executes tests automatically. Parametrized tests verify behaviour across multiple configurations (e.g., FFT sizes) without code duplication:

\begin{verbatim}
@pytest.mark.parametrize("n_channels", [64, 128, 256, 512, 1024])
def test_fft_linearity(n_channels):
    assert np.allclose(fft(signal1 + signal2), 
                       fft(signal1) + fft(signal2))
\end{verbatim}

Fixtures provide reusable test infrastructure, generating consistent configurations shared across tests.

The test infrastructure supports continuous integration, though not fully implemented in the research prototype.

\subsection{Performance Benchmarking}

Benchmark tests measure execution time and assert performance thresholds:

\begin{verbatim}
def test_correlation_throughput():
    start_time = time.time()
    run_correlator(config)
    assert time.time() - start_time < 5.0
\end{verbatim}

Performance assertions detect throughput regressions from inefficient code changes.

Validation ideally compares against trusted reference implementations like DiFX. The baseline testing compares against analytical solutions for specific scenarios. The methodology emphasises rigour balanced against resource constraints.
